\section{Trayectoria personal: de la App Store a emprender con agentes}

\subsection{Primera etapa: intuición de producto en el auge de la Internet móvil}

Ji Yichao (季逸超), durante la preparatoria, desarrolló un navegador de iOS de terceros y generó un ingreso de varios cientos de miles de dólares; esta experiencia forjó dos aprendizajes que siguieron siendo válidos a largo plazo: primero, los usuarios están dispuestos a pagar por «velocidad y tranquilidad»; segundo, la ventana de oportunidad tecnológica es efímera, y la velocidad misma es una ventaja competitiva. El modelo Buy Copy de aquella primera etapa era simple, pero entrenó la capacidad más difícil de adquirir: traducir directamente decisiones técnicas en flujo de efectivo.

\begin{importantbox}{Capacidades sólidas forjadas en el emprendimiento temprano}
\begin{itemize}
    \item Traducir las «quejas de los usuarios» en requisitos de ingeniería, en lugar de convertirlas en una visión abstracta
    \item Cuando el cómputo y el personal son limitados, priorizar el «punto de dolor de mayor frecuencia»
    \item Sustituir KPI complejos por métricas simples y verificables
\end{itemize}
\end{importantbox}

\subsection{Segunda etapa: de las necesidades de precarga a la investigación en PLN}

El detonante para pasar del negocio del navegador al PLN no fue un «interés académico», sino un problema de producto: cómo predecir el siguiente clic y precargar contenido con anticipación. Este problema es, en esencia, un problema de modelado de secuencias. Tras la publicación de Word2Vec, se contó con representaciones densas de texto, y el equipo de ingeniería finalmente tuvo la oportunidad de convertir reglas empíricas en una capacidad de modelo generalizable.

\begin{knowledgebox}{El costo real de la migración tecnológica}
Un punto que se repite en la entrevista: el cambio generacional de tecnología no le da al sistema anterior la oportunidad de una «actualización sin fricciones». De Word2Vec a LSTM, y luego a Transformer/BERT, cada ronda exigió reescribir parte del canal de datos, las especificaciones de etiquetado y los scripts de entrenamiento.
\end{knowledgebox}

\subsection{Tercera etapa: emprendimiento con grafos de conocimiento y análisis retrospectivo del fracaso}

El segundo emprendimiento se centró en la búsqueda semántica y la extracción abierta de información (Open IE), con el objetivo de superar el paradigma tradicional de los «diez enlaces azules». El equipo realizó una gran cantidad de trabajo de base: preentrenamiento de modelos, índices vectoriales, construcción de grafos, pero a nivel de producto no se logró un crecimiento sostenido. El análisis retrospectivo en la entrevista es muy directo: que la tecnología sea superior no implica que el negocio sea viable, sobre todo cuando no se tiene ventaja en el ecosistema, los canales ni la autorización de datos.

\begin{warningbox}{Tres manifestaciones concretas de llegar «demasiado pronto»}
\begin{itemize}
    \item Los puntos de interacción aún no maduraban: los dispositivos portátiles y las interfaces de voz no habían formado una distribución dominante
    \item Se subestimaron las barreras no técnicas: el suministro de datos y las relaciones de colaboración determinaban el límite superior
    \item El modelo de negocio oscilaba: el cambio entre 2C y 2B provocó una fractura en los objetivos organizacionales del equipo
\end{itemize}
\end{warningbox}

\subsection{Cuarta etapa: experiencia práctica con LLM dentro de una empresa y la decisión de volver a emprender}

Después de unirse a una empresa unicornio, Ji Yichao sistematizó, dentro del «sistema de rankings», una forma de pensar basada en benchmarks: cómo convertir necesidades en métricas comparables, y cómo alinear la asignación de recursos con los resultados de evaluación. Tras GPT-3, el paradigma de la industria cambió, y él volvió a ver una ventana para emprender. En comparación con el primer emprendimiento, esta vez se enfatizó más la «ingeniería de sistemas + circuito cerrado de producto», en lugar de una capacidad de modelo puntual.

\begin{importantbox}{De «hacer modelos» a «hacer sistemas»}
El cambio más digno de recordar en la entrevista es el cambio en la función objetivo: antes se enfocaba en las métricas del modelo; después, en la tasa de finalización confiable de la tarea (task completion reliability). Esto fue lo que determinó que Manus eligiera la ruta de sistema de «modelo + sandbox + herramientas».
\end{importantbox}

\subsection{Resumen de la sección}

La trayectoria de Ji Yichao no es una escalada lineal, sino una serie de «giros impulsados por problemas». Lo que en verdad es transferible no es una pila de modelos en particular, sino tres acciones: validar con rapidez, detener las pérdidas a tiempo, y convertir el fracaso en las condiciones de frontera de la siguiente etapa.
